\input{../tex-inputs/latex-leseansicht-vorspann}

\section[Gustav Schwarzkopf an Arthur Schnitzler, 12. 7. 1911]{L04263 Gustav Schwarzkopf an Arthur Schnitzler, 12. 7. 1911}
\nopagebreak\mylabel{L04263v}
\rehead{ }\normalsize\beginnumbering\briefempfaengerindex{Schnitzler, Arthur@\textsc{Schnitzler, Arthur}!zzzSchwarzkopf, Gustav@\emph{von Gustav Schwarzkopf}!1911-07-121@{12. 7. 1911}|(be}
\toendnotes[C]{\smallbreak\pagebreak[2]}
\correspDesc{Versand  durch Gustav Schwarzkopf am 12. 7. 1911 in Mautern in Steiermark
\newline{}Erhalt  durch Arthur Schnitzler im Zeitraum [13. 7. 1911 – 17. 7. 1911?] in Wien}\toendnotes[C]{\smallbreak}
\Standort{CUL, Schnitzler, B 96a.}
\physDesc{Bildpostkarte, 236 Zeichen
\newline{}Handschrift: schwarze Tinte, deutsche Kurrent
\newline{}Versand: Stempel: »\nobreak{}\oindex{Mautern in Steiermark@\textbf{Mautern in Steiermark}, \emph{Verwaltungsgebiet}|pwk}{[}Mautern{]} in
                                       Steiermark, {[}1{]}\textcolor{gray}{2}. VII. 11\nobreak{}«.  
\newline{}Schnitzler: mit Bleistift Vermerk: »\textsc{Schwarzkopf}« }\toendnotes[C]{\smallbreak}\pstart{}{\pb}Herrn\pend{}\pstart{}\textsc{D\textsuperscript{r} Arthur Schnitzler}\pend{}\pstart{}\textsc{Wien (Cottage)\oindex{Wien@\textbf{Wien}!XVIII., Währing@\textbf{XVIII., Währing}!Währinger Cottage@\textbf{Währinger Cottage}, \emph{Teil eines besiedelten Ortes}|pw}}\pend{}\pstart{}Sternwarteſtraße 71\oindex{Wien@\textbf{Wien}!XVIII., Währing@\textbf{XVIII., Währing}!Sternwartestraße 71@\textbf{Sternwartestraße 71}, \emph{Wohngebäude}|pw}\pend{}{\bigskip}
\pstart
           \noindent{}\centering{}{\pb}\textcolor{gray}{\textbf{Mautern, Steiermark\oindex{Mautern in Steiermark@\textbf{Mautern in Steiermark}, \emph{Verwaltungsgebiet}|pw}.}}\pend
           \vspace{1em}
\pstart
           \raggedleft{}{\pb}Mautern\oindex{Mautern in Steiermark@\textbf{Mautern in Steiermark}, \emph{Verwaltungsgebiet}|pw}{ }12. 7.\pend
           
\pstart
           \raggedleft{}Gaſthof \textsc{Klossner}\oindex{Gasthof Josef Kloßner@\textbf{Gasthof Josef Kloßner}, \emph{Gastgewerbegebäude}|pwv}\pend
           \vspace{0.5em}
\pstart
           Es iſt hier{ }ſehr hübſch, noch{ }ſehr{ }ſchwach beſucht und man hat die beſte Gelegenheit
               das Baden zu verlernen.\pend
           
\pstart
           Ihre Frau\pwindex{Schnitzler, Olga 17.\,1.\,1882 Wien – 13.\,1.\,1970 Lugano@\textsc{Schnitzler, Olga} (17.\,1.\,1882 Wien – 13.\,1.\,1970 Lugano), \emph{Schauspielerin, Sängerin}|pwv} und{\\[\baselineskip]}Sie herzlichſt{\\[\baselineskip]}grüßend{\\[\baselineskip]}Ihr{\\[\baselineskip]}\spacefill\mbox{Gustav.}\pend
           \leftskip=0em{}\selectlanguage{ngerman}\endnumbering\briefempfaengerindex{Schnitzler, Arthur@\textsc{Schnitzler, Arthur}!zzzSchwarzkopf, Gustav@\emph{von Gustav Schwarzkopf}!1911-07-121@{12. 7. 1911}|)be}\mylabel{L04263h}
\begin{anhang}
\end{anhang}\newcommand{\dateiname}{L04263}\newcommand{\titel}{Gustav Schwarzkopf an Arthur Schnitzler, 12. 7. 1911}\newcommand{\editorInnen}{Herausgegeben von Jahnke, SelmaMüller, Martin Anton}\input{../tex-inputs/latex-leseansicht-abspann}
